%==============================================================================
%  Chapter 4 — Findings
%
%  Organised around synthesis_v2.md (Pass 2 closed):
%    §1   Logic profile by section (orienting, with logic succession headline)
%    §2   Pattern A — Tension-Dissolution
%    §3   Pattern B — Coercive Governance Cluster
%    §4   Pattern C — State Social Equity Template
%    §5   Pattern D — Challenge Chapter Logic Succession Template
%    §6   Pattern E — STEP Circular Market Validation
%    §7   Pattern F — TRL as Cross-Pillar Legitimation Artefact
%    §8   STEP additional findings (state vocabulary as register,
%         Sovereignty Seal, hybrid instrument, mimetic isomorphism)
%    §9   Sustainability as adaptive state-policy entry vector
%         (candidate template not confirmed; four-mechanism inventory)
%    §10  Chapter summary
%
%  Keep the chapter descriptive — interpretation belongs in Chapter 5.
%  Entries are cited as [NNN] referring to Readthrough_Notes_v1.md;
%  patterns and coding rules are referenced from
%  Codebook_v2_Inductive.md (§7 patterns, §9 coding rules).
%==============================================================================
\chapter{Findings}
\label{ch:results}

\section*{Chapter overview}

This chapter presents the empirical findings of the readthrough of the
Pathfinder and STEP Scale Up sections of the EIC Work Programme 2026,
following the order in which the synthesis builds them. The chapter
opens with a logic profile across the corpus, then sets out the six
patterns confirmed in Pass~2 (Patterns A--F), then collects the
STEP-specific findings that do not themselves form patterns, and
closes with the sustainability-entry inventory, the candidate template
that the analysis did not confirm. The register throughout is
descriptive: the analytical interpretation of these findings belongs
in Chapter~\ref{ch:discussion}.

%==============================================================================
\section{Logic profile by section}
\label{sec:res-profile}

The corpus distributes its institutional logics unevenly across the
two instruments. The Pathfinder side opens in
\texttt{LOGIC\_SCIENCE} at the level of objectives, surrounds it with
\texttt{LOGIC\_STATE} in scope and impact framings, and routes its
governance through \texttt{LOGIC\_PROFESSION}. The STEP side is
dominated by \texttt{LOGIC\_STATE}, with \texttt{LOGIC\_MARKET}
genuinely operative only in the award criteria of Table~9 and a small
set of investment-mechanics passages. The profile in
\cref{tab:res-profile} reproduces \texttt{synthesis\_v2.md~\S1}; the
rest of the chapter unpacks the patterns that produce it.

\begin{table}[H]
\centering
\caption[Logic profile by section]{Dominant institutional logics by
section of the EIC Work Programme 2026 (Pathfinder + STEP Scale Up).
Source: \texttt{synthesis\_v2.md~\S1}, drawing on the Pass~2 codings
in \texttt{Readthrough\_Notes\_v1.md}.}
\label{tab:res-profile}
\begin{tabularx}{\textwidth}{@{}
  >{\raggedright\arraybackslash}p{0.26\textwidth}
  >{\raggedright\arraybackslash}p{0.27\textwidth}
  X @{}}
\toprule
\textbf{Section} & \textbf{Dominant logic(s)} & \textbf{Notes} \\
\midrule
Overview / Introduction
  & \texttt{STATE} + \texttt{MARKET}
  & Pattern~A signature at the Strategic Goals level; constructed-argument \texttt{LOGIC\_TENSION} codings withdrawn in Pass~2 per coding rule~1. \\
Pathfinder Open --- Key features
  & \texttt{STATE} + \texttt{PROFESSION}
  & Programme Manager authority cluster (Pattern~B); rule~3 brings [024] into \texttt{ISO\_COERCIVE} + \texttt{LOGIC\_STATE}. \\
Pathfinder Open --- Evaluation
  & \texttt{SCIENCE} + \texttt{STATE}
  & Science criteria colonised by state and market vocabulary ([056], [057]); Pattern~A holds tension below the surface. \\
Pathfinder Challenges --- Governance
  & \texttt{PROFESSION} + \texttt{STATE}
  & Portfolio architecture; Programme Manager cluster culminates in [061] (Pattern~A and Pattern~B intersect). \\
II.2.1 Advanced Materials
  & Pattern~D template
  & \texttt{STATE} $\to$ \texttt{SCIENCE} $\to$ \texttt{PROFESSION} $\to$ \texttt{STATE} succession identified here first. \\
II.2.2 Healthy Ageing
  & Pattern~D template
  & Sustainability-candidate universality fails here (no environmental language present). \\
II.2.3 DeepRAP
  & Pattern~D template
  & \texttt{LOGIC\_STATE} enters via EU AI Act compliance ([129]); sustainability appears obliquely as compute-constraint efficiency proxy ([125], [129]). \\
STEP --- Access \& eligibility
  & \texttt{STATE}
  & Market vocabulary as state legitimation register (disciplined by coding rule~7). \\
STEP --- Award criteria (Table~9)
  & \texttt{MARKET} + \texttt{STATE} + \texttt{PROFESSION}
  & Clearest genuine \texttt{LOGIC\_MARKET} in the corpus ([169], [172], [173]); Pattern~E anchored at [140] + [176] with the 20--33\% anchor position at [145]. \\
\bottomrule
\end{tabularx}
\end{table}

\paragraph{Logic succession across the innovation pipeline.}
Reading the profile from top to bottom is reading the EIC's
innovation pipeline from low TRL to high. As the technology matures
from proof-of-principle to scale-up, the dominant institutional logic
shifts. \texttt{LOGIC\_SCIENCE} organises the Pathfinder Open
objectives at TRL~1--4. \texttt{LOGIC\_STATE} progressively organises
the impact framings of the Pathfinder Challenges.
\texttt{LOGIC\_STATE} dominates the STEP architecture at
commercialisation stage, with \texttt{LOGIC\_MARKET} present mainly
as a constrained evaluative pole in Table~9. The succession is not a
smooth transition. It is punctuated by embedded market requirements
in science-phase Challenges ([102], [109]) and by science criteria in
the STEP award structure ([169]). These cross-logic embeddings are
where the patterns of the rest of the chapter live.

\paragraph{Intra-chapter reproduction.}
A second, finer-grained reproduction of the same succession is
observable inside each Challenge chapter (see
\cref{sec:res-pattern-d}, Pattern~D). Scope and background frame the
Challenge in \texttt{LOGIC\_STATE} terms; the Specific Objectives
shift to \texttt{LOGIC\_SCIENCE}; the Portfolio Approach routes
through \texttt{LOGIC\_PROFESSION}; the Expected Impact returns to
\texttt{LOGIC\_STATE}. The macro-level succession across the pipeline
is therefore reproduced at chapter scale across substantively
different technology domains. The meso-level reading of this
reproduction is taken up in \Cref{ch:discussion}.

%==============================================================================
\section{Pattern A --- Tension-Dissolution}
\label{sec:res-pattern-a}

Pattern~A names a recurring textual move: where two institutional
logics could be presented as in conflict, the text instead presents
them as naturally complementary, joining them with \enquote{and} or
listing them as parallel sub-criteria. The tension is structurally
present, but the text does not acknowledge it. The conflicts the
analysis is most interested in (open science against
intellectual-property exploitation, scientific curiosity against
market readiness, \enquote{any field} against \enquote{named
technology domains}) are exactly the ones the text passes over in
silence.

\paragraph{Distinction from \texttt{LOGIC\_TENSION}.}
The \texttt{LOGIC\_TENSION} code applies only where the conflict is
legible in the text itself, through hedging, explicit acknowledgement
of competing priorities, or structural segmentation that visibly
contains the tension. Pass~2 introduced this as coding rule~1 (the
surface-visibility standard for \texttt{LOGIC\_TENSION}), and
Pattern~A is its analytical counterpart: it names the cases where
conflict is structurally present but rhetorically suppressed. Several
Pass~1 codings of \texttt{LOGIC\_TENSION} were withdrawn in Pass~2
once rule~1 was formalised, including [003], [008], and [019]. In
each case the underlying configuration is a Pattern~A signature
rather than a surface-visible tension.

\paragraph{Instances.}
The pattern is observable across the corpus, but it is most
concentrated in three locations.

\begin{itemize}[leftmargin=*]
  \item \textbf{Overview and Strategic Goals.} The Strategic Goals
    paragraph at [003] presents science, market, and state logics as
    co-implicated in a single \enquote{research base across the EU},
    quietly reframing what counts as valuable science. The Overview
    paragraph at [008] arranges \enquote{any field}
    (\texttt{LOGIC\_SCIENCE}) and \enquote{named technology domains}
    (\texttt{LOGIC\_STATE}) as a continuous arc of instrument design
    rather than as alternatives.
  \item \textbf{Pathfinder Open evaluation criteria.} The Excellence
    criterion at [056] folds market vocabulary
    (\enquote{high-risk/high-gain}, \enquote{radically new
    technology}) into a science-quality criterion without
    acknowledging that the two pull in different directions. The
    Impact criterion at [057] lists market success and societal or
    environmental impact as equivalent sub-criteria. The gatekeepers
    passage at [044] presents three governance criteria (what the
    instrument produces, how it works, and who it admits) as
    complementary rather than as drawn from different logics.
  \item \textbf{Challenge Specific Objectives.} [102] self-identifies
    as a Pattern~A instance: \enquote{potential markets and
    associated impacts} are appended to science objectives via
    \enquote{also}, subordinate but mandatory. [109] performs the
    same move at the Healthy Ageing Specific Objectives, listing
    science proof-of-concept and market-exploitation deliverables as
    co-equal requirements. The tension surfaces only when [108],
    which acknowledges the low translational success rate in ageing
    research, is read against [109].
\end{itemize}

\paragraph{The most consequential instance.}
[082] embeds state-defined Challenge relevance as the primary
sub-criterion of scientific excellence. Applicants must frame their
science to fit a state-defined Challenge in order to receive funding,
and the text presents this as a quality-of-science criterion rather
than as a state-organising one. This is Pattern~A operating at the
gate of access: it shapes how science is allowed to be science in the
Challenge calls.

\paragraph{Pattern~A as a textual strategy.}
The instances above share a structural feature: logics that conflict
in practice are presented as if they did not. The strategy performs
institutional coherence rather than arguing for it; it is silent
about exactly the conflicts the readthrough surfaces. Whether this
silence is best read as a cognitive-legitimacy strategy presenting
the instrument as obviously coherent, as a managed compromise between
mandates that cannot otherwise be reconciled in a single text, or as
something else, is the interpretive question taken up in
\Cref{ch:discussion}.

%==============================================================================
\section{Pattern B --- Coercive governance cluster}
\label{sec:res-pattern-b}
\todoinline{Programme Manager discretionary authority over capital
continuity, roadmap direction, portfolio placement, and
inter-consortium collaboration, in tension with the autonomy
protections stated at [021]. Walk the intensity gradient from
(i)~stated autonomy ([021]) through (ii)~milestone oversight ([023])
to (iii)~suspension or termination authority ([024]), now coded
\texttt{ISO\_COERCIVE} + \texttt{LOGIC\_STATE} under rule~3
(contractual discretionary coercion).}

\todoinline{Continuing the gradient: (iv)~portfolio-placement direction
([025], \enquote{may be requested to join}); (v)~proactive
roadmap-setting and steering ([026], [027]); the frozen-notes
consequence is that refined rule~3 also implicates [026] for hard
\texttt{ISO\_COERCIVE}, logged at v2-006 but not retroactively
applied; (vi)~dual-responsiveness mandate without resolution ([061]);
(vii)~soft collaboration directives ([115], \enquote{encouraged to
collaborate}). The cluster forms an intensity gradient from hard
contractual sanction to soft directive language. The gradient is a
property of the cluster configuration, not of any single passage,
and is therefore recorded as Pattern~B (observation) rather than
promoted to a code.}

\todoinline{Substantive finding for \Cref{ch:discussion}: a hired
technocrat, neither a peer scientist nor an elected official, holds
decisive authority over capital continuity through legally neutral
grant-mechanics language. The gap between [021]'s stated autonomy and
[024] / [026]'s operative control is the descriptive backbone.}

%==============================================================================
\section{Pattern C --- State social-equity template}
\label{sec:res-pattern-c}
\todoinline{Near-identical passages embedding gender balance, female
researcher empowerment, and underrepresented-groups mandates recur
across instruments and sections with minimal variation. Anchor the
section on the [047] $\approx$ [063] near-identical pair (Pathfinder
Open team composition $\approx$ Pathfinder Challenges team
composition) as the clearest evidence of a template. Wider cluster:
[001], [017], [047], [063], [112], [175]. Tiebreaker reinforcement at
[059]: gender balance ranked above SME count and Member State
geographic representation in the tiebreaker architecture.}

\todoinline{Codebook-level consequence: coding rule~4 disciplines
\texttt{NORM\_DEI} / \texttt{LOGIC\_STATE} co-coding so that
state-mandated equity templates route through \texttt{LOGIC\_STATE};
\texttt{NORM\_DEI} is reserved for credible-plan organisational
standards with evaluation weight. Frozen-notes consequence: [001]'s
Pass~1 \texttt{LOGIC\_TENSION} coding is treated for Ch~4 purposes as
\texttt{LOGIC\_STATE} only (rule~1; logged at v2-011, not
retroactively applied to v1 notes).}

\todoinline{Descriptive ground: textual near-identity across
instruments indicates standardised insertion at the EU legislative
layer (Horizon Europe Regulation gender mainstreaming) rather than
EIC-specific programmatic commitment. Interpretive question deferred
to \Cref{ch:discussion}: why state-equity priorities are inserted
into an innovation funding instrument and ranked above market and
geographic priorities in the tiebreaker.}

%==============================================================================
\section{Pattern D --- Challenge-chapter logic succession template}
\label{sec:res-pattern-d}
\todoinline{Each Challenge chapter reproduces a consistent
intra-chapter logic succession across its sub-sections. Scope /
Background $\to$ \texttt{LOGIC\_STATE} (state-defined strategic
relevance, sectoral framing); Specific Objectives $\to$
\texttt{LOGIC\_SCIENCE} (methodology, capability advancement);
Portfolio Approach $\to$ \texttt{LOGIC\_PROFESSION} (Programme
Manager direction, portfolio composition); Expected Outcomes /
Impacts $\to$ \texttt{LOGIC\_STATE} (societal, environmental,
sovereignty goals).}

\todoinline{Instances: [095]--[116] (II.2.1 + II.2.2); [117]--[129]
(II.2.3). The template was first identified at II.2.1 and confirmed
across II.2.2 and II.2.3, three substantively different technology
domains (advanced materials for energy harvesting; biotechnology for
healthy ageing; AI / deep reinforcement and adaptive planning)
reproducing the same logic configuration. Codebook-level result:
pure analytical observation, with no new code, no flag-row
recoding, and no coding rule. Underlying passage codings are
individually correct; Pattern~D is the configuration property.}

\todoinline{Substantive finding for \Cref{ch:discussion}: the state's
institutional reproduction operates at the chapter-template level,
independent of technological content. This is a meso-level
isomorphism that the passage-level \texttt{ISO\_*} codes do not
capture. Pattern~D also enriches the logic-succession finding from
\cref{sec:res-profile}: the macro-level succession across the
pipeline is reproduced as an intra-chapter succession in every
Challenge.}

%==============================================================================
\section{Pattern E --- STEP circular market validation}
\label{sec:res-pattern-e}
\todoinline{STEP simultaneously requires two things of every
applicant. First, \emph{demonstrated market interest}: a
non-legally-binding pre-commitment from a qualified investor for at
least 20\% of the target round ([141]). Second, \emph{demonstrated
market failure}: the amount cannot be fully financed by private
investors without EIC participation ([176]).}

\todoinline{These two requirements are only simultaneously satisfiable
by using EIC's prospective 20--33\% anchor participation ([145]) as
the lever to obtain the private pre-commitment. The private investor
commits because EIC will co-invest; the EIC requires the commitment
as proof of market interest. The market signal the state requires is
constructed by the state's own potential participation. At this
scale of anchor position, private investors are not independently
validating the company; they are co-investing around a dominant
state anchor.}

\todoinline{Codebook-level result: pure analytical observation plus
coding rule~2 (the STEP catalytic-investment co-coding rule).
Rule~2 disciplines \texttt{LOGIC\_STATE} + \texttt{LOGIC\_MARKET}
co-coding in catalytic-investment passages: code \texttt{LOGIC\_STATE};
co-code \texttt{LOGIC\_MARKET} only where market mechanics are
operationally surface-visible in the passage. The structural
circularity itself is a relational property across the configuration
[132], [140], [141], [145], [176], not of any single passage, and is
therefore recorded as Pattern~E (observation), not as a
\texttt{LOGIC\_TENSION} coding.}

\todoinline{Description targets for the prose: walk the four entries
in order ([132] catalytic mechanism $\to$ [140] / [141]
market-interest eligibility criterion $\to$ [145] 20--33\% anchor
$\to$ [176] market-failure award criterion). The interpretive
question (fundamental contradiction or deliberate developmental-state
design) is deferred to \Cref{ch:discussion}.}

%==============================================================================
\section{Pattern F --- TRL as cross-pillar legitimation artefact}
\label{sec:res-pattern-f}
\todoinline{TRL (Technology Readiness Level) operates simultaneously
across two of Scott's pillars within the same passages. As
\texttt{REG\_TRL} (regulative pillar), TRL performs the gate
function: eligibility threshold, scope boundary, budget allocation,
and instrument routing. As \texttt{COG\_LINEAR} (cognitive pillar),
TRL is the taken-for-granted developmental ladder, borrowing
measurement legitimacy from its NASA / US-defence origin.}

\todoinline{The two codes are not in tension; they co-code in the
same passages. The cognitive pillar provides scientific-register
surface cover; the regulative pillar performs the gatekeeping
beneath. The co-operation depoliticises political decisions about
where in the innovation pipeline public money flows. Cited entries:
[012], [044], [046]; the configuration recurs across Pathfinder
evaluation passages.}

\todoinline{Codebook-level result: pure analytical observation. v1
definitions of \texttt{REG\_TRL} and \texttt{COG\_LINEAR} are
unchanged; v1 already permits co-coding. The disguise is a
relational property of how two correctly-applied codes interact,
not a coding ambiguity to resolve. Collapsing the codes would erase
the contrast that makes the legitimation work readable.}

\todoinline{Triple-load finding (Pass~2 addition): the TRL artefact
does legitimation work along three axes simultaneously, cognitive
(scientific-measurement register), regulative (eligibility gate),
and mimetic (NASA / US-defence provenance; see also
\cref{sec:res-step-mimesis} on the DARPA / Temasek mimetic
isomorphism finding). Description targets for the prose: name each
axis with a passage anchor; reserve the cross-pillar interpretation
(Scott + DiMaggio \& Powell + Suchman intersecting in a single
artefact) for \Cref{ch:discussion}.}

%==============================================================================
\section{STEP additional findings}
\label{sec:res-step-additional}

\todoinline{Introductory paragraph: STEP exhibits three further
findings that complement Pattern~E without themselves forming
patterns at the same level: the state's absorption of market
vocabulary as its own legitimation register, the Sovereignty Seal
as a portable legitimacy artefact, and the explicit mimetic
benchmarking against DARPA and Temasek. The section closes with a
hybrid-instrument table summarising how STEP combines private-equity
mechanics with state tender architecture.}

\subsection{State absorbing market vocabulary as legitimation
register}
\label{sec:res-step-vocabulary}
\todoinline{Across the STEP section (anchored at [130], developed
across [132]--[180]), \texttt{LOGIC\_STATE} absorbs and deploys
market vocabulary (\enquote{market gap}, \enquote{Venture Debt},
\enquote{equity}, \enquote{deal flow}) as its own legitimation
register. The vocabulary signals the state justifying its
intervention through market-failure language, not market logic
governing the passage.}

\todoinline{Pass~2 codifies this distinction via coding rule~7 (the
vocabulary-vs-organising standard): a logic is coded only where it
organises the passage; vocabulary-level signals in procedural,
administrative, or budgetary passages do not in themselves warrant
primary coding. Cite the [012] precedent
(state-as-administrative-funder below threshold) and the [014]
precedent (VC vocabulary in budgetary sentences below threshold).}

\todoinline{Exceptions where \texttt{LOGIC\_MARKET} genuinely
organises: Table~9 market-opportunity criterion ([172]),
business-model criterion ([173]), and innovation criterion's cost /
performance dimension ([169]). These are the passages where market
logic actually governs the evaluative frame.}

\subsection{The Sovereignty Seal as field-bridging legitimacy
artefact}
\label{sec:res-step-seal}
\todoinline{The Sovereignty (STEP) Seal is a state-issued credential
that travels across institutional fields. It unlocks access to
InvestEU Venture Debt ([152], [161]), enables access to Business
Acceleration Services, and, crucially, can be awarded even on NO~GO
decisions where merit criteria are met but budget is unavailable
([166]). The decoupling from the investment decision reveals the
Seal as a pure legitimacy instrument: the state certifies strategic
value independently of whether it commits capital. Description
targets: describe the Seal mechanism, the NO~GO award path, and the
inter-institutional reach (public EU programmes and private
investors). The interpretive reading (Suchman pragmatic and moral
legitimacy via a designed artefact) is for \Cref{ch:discussion}.}

\subsection{STEP as hybrid instrument}
\label{sec:res-step-hybrid}
\todoinline{STEP deploys private-equity mechanics through state
tender infrastructure. Neither logic fully governs the instrument as
a whole; both are operative in different dimensions. Reproduce the
hybrid-instrument table from \texttt{synthesis\_v2.md~\S2.9}: market
mechanics column (equity-only, case-by-case investment negotiation,
KYC/AML investor screening, 3--5$\times$ leverage requirement,
Investment Committee decision) against state access-architecture
column (quarterly evaluation batches, Commission Award Decision,
jury of maximum six members, 50-page business plan + FTO analysis,
Funding and Tenders Portal submission). The descriptive point:
market and state operate side by side in the same instrument
without integration, and the architecture is what coding rule~2
disciplines at the passage level.}

\subsection{Mimetic isomorphism against DARPA, Temasek and Chinese
GGFs}
\label{sec:res-step-mimesis}
\todoinline{The Work Programme explicitly benchmarks STEP against US
(DARPA) and Asian (Temasek, Chinese Government Guidance Funds)
models. Document the citation; characterise the adoption as
selective: STEP reproduces the equity / catalytic-investment
mechanics while embedding them in EU procurement architecture
(Funding and Tenders Portal, quarterly batches, Commission Award
Decision, jury process, 50-page business plan + FTO analysis). The
result is a hybrid that does not fully replicate the models it
cites. Cross-reference \cref{sec:res-pattern-f}: the TRL framework
itself carries DARPA / US-defence provenance, so mimetic isomorphism
operates across Pathfinder (TRL ladder) and STEP (equity instrument)
alike, not only at the instrument-design level of STEP.}

%==============================================================================
\section{Sustainability as an adaptive state-policy entry vector}
\label{sec:res-sustainability}
\todoinline{Pass~1 raised the hypothesis (at [096]) that
sustainability framing might recur across Challenge chapters as a
standardised state-logic insert, analogous to Pattern~C (the State
Social Equity Template). Pass~2 tested the universality of this
candidate template and did not confirm it: at II.2.2 (Biotechnology
for Healthy Ageing) there is no sustainability or environmental
language present. Sustainability is therefore
subject-matter-dependent, not a universal Challenge insert.}

\todoinline{The substantive finding from the test is preserved as a
four-mechanism inventory: when sustainability does appear, it enters
via one of four distinct mechanisms. (1)~\emph{Commission
Communication policy-citation} ([099]): sustainability flows through
institutional inheritance from a cited Commission Communication
(RePowerEU, Green Deal). (2)~\emph{Explicit
environmental-consequence framing} ([096]): sustainability appears
as descriptive context establishing the problem the Challenge
addresses (IoT energy consumption, battery waste).
(3)~\emph{Sector-definition} (around [117]--[120]): sustainability
is constitutive of the technology field itself (clean and
resource-efficient technologies as a STEP critical-area category).
(4)~\emph{Technical efficiency-proxy} ([125] / [129]): sustainability
is absorbed into methodology requirements via constrained-resource
specifications (compute-constraint requirements in DeepRAP), without
explicit environmental framing or policy citation.}

\todoinline{Analytical contrast with Pattern~C (deferred to
\Cref{ch:discussion}): Pattern~C operates as a universal insert
reproduced in near-identical wording regardless of subject matter;
sustainability operates as an adaptive entry vector that adapts its
expression to Challenge content. The institutional architecture
therefore carries a small set of state priorities universally
(gender mainstreaming) while letting other state priorities
(sustainability) enter through subject-matter-appropriate channels.}

%==============================================================================
\section{Chapter summary}
\label{sec:res-summary}
\todoinline{Short summary closing the chapter. Touch in order: the
logic profile and the Science $\to$ Market $\to$ State succession
across TRL (\cref{sec:res-profile}); the six confirmed patterns A--F
(\cref{sec:res-pattern-a,sec:res-pattern-b,sec:res-pattern-c,sec:res-pattern-d,sec:res-pattern-e,sec:res-pattern-f});
the STEP additional findings (\cref{sec:res-step-additional}); the
sustainability inventory and the negative-universality result
(\cref{sec:res-sustainability}). Close with the descriptive-to-
interpretive handover: the patterns above are presented in their
empirical shape; their theoretical implications, including
Pattern~E's developmental-state reading, Pattern~F's triple-load
across Scott's pillars, and Pattern~C's social-equity placement
above market and geographic priorities, are for
\Cref{ch:discussion}.}
